% Lucky passes paper: CODEX version (Journal of Computing, Data, and Exploration, U of T Mississauga).
% Template: https://www.overleaf.com/read/wnzmspydjxrv (main.tex + pepadun.sty + pepadun.bib), read 2026-11-25.
% Build: uv run python scripts/build_overleaf.py --venue codex -> paper/codex_overleaf.zip. On Overleaf, open the
% template, Menu -> Copy Project, then upload the zip's files into the copy (replace main.tex; keep pepadun.sty and
% pepadun.bib, which this file needs). Compiler: pdfLaTeX; bibliography: biber (the style loads biblatex, IEEE).
% Same section files as the ACL version (paper/main.tex); only the frame differs. Numbers checked by check_numbers.
%
% Journal rules this file follows (research-progress.md, "Venue: CODEX"):
% - double-blind: the title page has the title only; no author, tool name or repository link;
% - title of at most 10 words; abstract 150-250 words (submissions page; the template says 200-300, so aim for
%   200-250); 3-5 keywords, alphabetical, separated by semicolons; no symbols or equations in title or abstract;
% - sections: Introduction, Research Methodology, Results and Discussion, Conclusions (template headings);
% - figures and tables inside the text where first discussed, numbered, referred to by number ("Table 1"), never
%   by position ("below"); tables with horizontal rules only (booktabs);
% - IEEE numeric citations; DOIs or URLs where available;
% - any LLM use disclosed in the acknowledgements.
\documentclass[11pt]{article}
\usepackage{pepadun}
\addbibresource{references.bib}

% The section files use natbib commands (ACL); map them to biblatex's IEEE style.
\newcommand{\citep}[1]{\cite{#1}}
\newcommand{\citet}[1]{\textcite{#1}}

\newcommand{\toolname}{an open-source RAG testing tool (name withheld for review)}
\newcommand{\repolink}{[anonymised repository link]}
\newcommand{\aliwrites}[1]{\noindent\fbox{\parbox{\linewidth}{\small\textbf{[Ali writes this section.]}\\ #1}}}

% At most 10 words. Ali picks it (paper/FRAMING.md).
\title{[Title: at most ten words; Ali picks]}
\author{}
\date{}

\begin{document}
\maketitle

\begin{abstract}
\aliwrites{Between 150 and 250 words, in English, with no symbols or equations. One paragraph each for: the problem; the method; % nocheck (word-count guidance, not a result)
the main result (one number with its interval and n, from results/FINDINGS.md); what a practitioner should do.
What do we add beyond Jonah et al.\ (2026)?}
\end{abstract}

\keywords{[3 to 5, alphabetical; e.g. evaluation; large language models; parametric knowledge; retrieval-augmented generation]}

\section{INTRODUCTION}
\input{sections/introduction}
\subsection{Related work}
\input{sections/related}

\section{RESEARCH METHODOLOGY}
\label{sec:method}
\input{sections/method}
\subsection{Reproducibility}
\input{sections/reproducibility}

\section{RESULTS AND DISCUSSION}
\label{sec:results}
\input{sections/results}
\subsection{Discussion}
\input{sections/discussion}
\subsection{Limitations}
\input{sections/limitations}

\section{CONCLUSIONS}
\input{sections/conclusion}

\begin{acknowledgments}
\input{sections/acknowledgements}
\end{acknowledgments}

\printbibliography[title={REFERENCES}, heading=bibintoc]

\end{document}
